\chapter{一変数の微分}
\label{ch:tb-differential}

\section{出題された二つの対数微分}

$u\ne0$ では $(\log\abs u)'=u'/u$ だから
\[
\frac{\d}{\d x}\log\abs{\tan x}
=\frac{\sec^2x}{\tan x}
=\frac1{\sin x\cos x}.
\]
また変数が指数にあれば、正の範囲で指数・対数へ戻す。例えば
\[
x^{\sin x}=\exp(\sin x\log x),
\qquad
\frac{\d}{\d x}x^{\sin x}
=x^{\sin x}\left(\cos x\log x+\frac{\sin x}{x}\right).
\]

\section{最大値・最小値}

\begin{theorem}{最大値・最小値の存在}{tb-extreme-value}
コンパクト集合上の連続関数は最大値と最小値を取る。
特に $\R^n$ の有界閉集合はコンパクトである。
\end{theorem}

\begin{proposition}{無限遠で0へ収束する関数}{tb-extreme-at-infinity}
連続関数 $f:\R\to\R$ が $\abs{x}\to\infty$ で $f(x)\to0$ なら、
$f$ は最大値または最小値を持つ。
\end{proposition}

\begin{proof}
$f\equiv0$ ならよい。$f(p)\ne0$ を取り、$\varepsilon=\abs{f(p)}/2$ とする。
ある $R>\abs p$ で
\[
\abs x>R\Longrightarrow\abs{f(x)}<\varepsilon
\]
となる。$[-R,R]$ 上で取る最大・最小を $M,m$ とする。
$f(p)>0$ なら $M\ge2\varepsilon$ なので $M$ は $\R$ 全体の最大値、
$f(p)<0$ なら同様に $m$ が全体の最小値である。
\end{proof}

具体的な最大・最小は、定義域、端点、$f'=0$ となる点、無限遠での極限を調べ、
$f'$ の符号で比較する。制約に平方根があれば
$u=\sqrt x$, $v=\sqrt y$ と置いて一変数へ落とす。

\section{一様収束と誤差関数}

$f_N:I\to\R$ が $f$ に一様収束するとは
\[
\sup_{x\in I}\abs{f_N(x)-f(x)}\longrightarrow0
\]
となることをいう。

有界閉区間 $[a,b]$ 上で連続な $f_N$ が $f$ に一様収束すれば
\[
\int_a^bf_N(x)\,\d x\longrightarrow\int_a^bf(x)\,\d x.
\]
従って関数項級数の部分和が一様収束すれば項別積分できる。

2026年度第1期では、指数関数についての次の有限和と誤差だけを使う。
ある $0<\theta<1$ に対して
\[
e^t=\sum_{n=0}^{N}\frac{t^n}{n!}
+\frac{e^{\theta t}t^{N+1}}{(N+1)!}.
\]
$t=-x^2$ とし、有界閉区間上で $\abs x\le M$ と取れば
\[
\left|e^{-x^2}-\sum_{n=0}^{N}\frac{(-1)^nx^{2n}}{n!}\right|
\le\frac{(M^2)^{N+1}}{(N+1)!}\longrightarrow0.
\]
これで必要な一様収束が直接示せる。項別積分により
\[
\int_0^xe^{-t^2}\,\d t
=\sum_{n=0}^{\infty}\frac{(-1)^n}{n!(2n+1)}x^{2n+1}.
\]
$x=1$ で $N$ 次まで使ったときは、積分後の誤差も
\[
\left|\int_0^1e^{-t^2}\,\d t
-\sum_{n=0}^{N}\frac{(-1)^n}{n!(2n+1)}\right|
\le\frac1{(N+1)!(2N+3)}
\]
と評価できる。

\section{加法的関数}

\begin{theorem}{連続な Cauchy の関数方程式}{tb-cauchy-equation}
$f:\R\to\R$ が $f(x+y)=f(x)+f(y)$ を満たし、連続なら、
ある $c\in\R$ により $f(x)=cx$ と書ける。
\end{theorem}

\begin{proof}
加法性から $f(0)=0$, $f(-x)=-f(x)$ を得る。整数 $m$ と正整数 $n$ について
$f(mx)=mf(x)$, $f(x/n)=f(x)/n$ だから、有理数 $r$ では $f(rx)=rf(x)$ である。
$c=f(1)$ とし、$x$ に収束する有理数列 $r_k$ を取る。連続性から
\[
f(x)=\lim_{k\to\infty}f(r_k)=\lim_{k\to\infty}cr_k=cx.
\]
\end{proof}
